\documentclass{article}
\usepackage{amsmath}
\usepackage{amssymb}
\usepackage{geometry}
\usepackage{tikz}
\usetikzlibrary{3d,calc}
\usepackage{tikz-3dplot}

\geometry{margin=1in}

\title{Configuration and Modeling of the KEEP System}
\author{Kite Electrical Energy Production (KEEP) Project}
\date{}

\begin{document}

\maketitle

\section{Configuration of the system and physical parameters}

\subsection{Configuration}

The system consists of an oscillating rigid arm of length $a$ and a kite $K$ connected to the tip of the arm $A$ via a tether of length $\ell$. The arm is pivoted at the origin $O$ and rotates in the horizontal $xy$-plane. 

\tdplotsetmaincoords{70}{-15}
\begin{figure}[htbp]
\centering
\begin{tikzpicture}[tdplot_main_coords, scale=2.2]

% ==== ORIGIN AND AXES ====
\coordinate (O) at (0,0,0);
\node[circle,fill,inner sep=1pt] at (O) {};
\node[below left] at (O) {O};

\draw[->] (0,0,0) -- (3,0,0) node[anchor=west]{$x$};
\draw[->] (0,0,0) -- (0,1.5,0) node[anchor=south]{$y$};
\draw[->] (0,0,0) -- (0,0,2) node[anchor=west]{$z$};

% ==== POINT A IN XY-PLANE ====
\def\a{2}
\def\alphaA{20}
\pgfmathsetmacro{\xA}{\a*cos(\alphaA)}
\pgfmathsetmacro{\yA}{\a*sin(\alphaA)}
\coordinate (A) at (\xA, \yA, 0);
\node[circle,fill,inner sep=1pt] at (A) {};
\node[below] at (A) {A};

% Projections of A
\coordinate (Ax) at (\xA, 0, 0);
\coordinate (Ay) at (0, \yA, 0);

% Lines from O
\draw[thick] (O) -- (A) node[midway, sloped, above] {$a$};
\draw[dotted] (O) -- (Ax);
\draw[dotted] (O) -- (Ay);

% Resisting torque indication at O
\draw[->, purple, thick] (0.4, 0, 0) arc (0:120:0.4) node[midway, above right] {$\Gamma_r$};

% Angle alpha
\tdplotdefinepoints(0,0,0)(1,0,0)(\xA,\yA,0)
\tdplotdrawpolytopearc[->]{1.5}{anchor=east}{$\alpha$}

% ==== POINT K ====
\def\l{2}
\def\thetaZero{70}
\def\DeltaTheta{5}
\def\DeltaPhi{30}
\def\tauK{90}

\pgfmathdeclarefunction{thetaFunc}{1}{\pgfmathparse{\thetaZero + \DeltaTheta * sin(2*#1)}}
\pgfmathdeclarefunction{phiFunc}{1}  {\pgfmathparse{\DeltaPhi * sin(#1)}}
\pgfmathdeclarefunction{xTauFunc}{1} {\pgfmathparse{sin(thetaFunc(#1)) * cos(phiFunc(#1))}}
\pgfmathdeclarefunction{yTauFunc}{1} {\pgfmathparse{sin(thetaFunc(#1)) * sin(phiFunc(#1))}}
\pgfmathdeclarefunction{zTauFunc}{1} {\pgfmathparse{cos(thetaFunc(#1))}}

\pgfmathsetmacro{\dotprod}{\xA * xTauFunc(\tauK) + \yA * yTauFunc(\tauK)}
\pgfmathsetmacro{\rK}{\dotprod + sqrt(\dotprod * \dotprod + \l^2 - \a^2)}

\pgfmathsetmacro{\xK}{\rK * xTauFunc(\tauK)}
\pgfmathsetmacro{\yK}{\rK * yTauFunc(\tauK)}
\pgfmathsetmacro{\zK}{\rK * zTauFunc(\tauK)}

\coordinate (K) at (\xK, \yK, \zK);
\coordinate (K_xy) at (\xK, \yK, 0);

% Helper points from A
\pgfmathsetmacro{\xAel}{\xA}
\pgfmathsetmacro{\yAel}{\yA}
\pgfmathsetmacro{\zAel}{1.5}
\coordinate (A_el) at (\xAel, \yAel, \zAel);

\pgfmathsetmacro{\xAext}{1.5*\xA}
\pgfmathsetmacro{\yAext}{1.5*\yA}
\coordinate (A_ext) at (\xAext, \yAext, 0);

% Draw K and lines
\node[circle,fill,inner sep=1pt] at (K) {};
\node[above] at (K) {K};

\draw (A) -- (K) node[midway, sloped, below] {$\ell$};
\draw[dotted] (O) -- (K) node[pos = 0.35, sloped, above] {$r$};

\draw[dotted] (A) -- (A_ext);
\draw[dotted] (K) -- (K_xy) -- (A);
\draw[dotted] (A) -- (A_el);

\newcommand{\KCoord}[2]{({#1 * xTauFunc(#2)}, {#1 * yTauFunc(#2)}, {#1 * zTauFunc(#2)})}

% Lemniscate Manifold
\draw[domain = 0:360, samples = 40, smooth, dashed, variable = \t]
    plot (\KCoord{\rK}{\t});
\draw[->, domain = 0:\tauK, samples = 20, smooth, pink, thick, variable = \t]
    plot (\KCoord{\rK}{\t});

\node[below, pink] at \KCoord{\rK}{\tauK/2} {$\tau$};

\end{tikzpicture}
\caption{Three-dimensional static configuration of the oscillating arm and the kite constrained to a spherical lemniscate manifold.}
\label{fig:static_config}
\end{figure}

\subsubsection{Coordinate Systems and Position Vectors}
We define a direct orthonormal reference frame $(\vec{e}_x, \vec{e}_y, \vec{e}_z)$ centered at $O$.
The position of the arm's tip $A$ is restricted to the horizontal $xy$-plane at a fixed distance $a$ from the origin:
\begin{equation}
\vec{x}_A(\alpha) = \begin{pmatrix} a \cos\alpha \\ a \sin\alpha \\ 0 \end{pmatrix}
\end{equation}
where $\alpha \in [0, 2\pi]$ represents the angle of the arm relative to the $x$-axis.

The kite $K$ is located at a distance $r$ from the origin along the direction vector $\vec{u}(\tau)$, which represents a lemniscate manifold (figure-eight trajectory) on the unit sphere:
\begin{equation}
\vec{x}_K(\tau, r) = r \vec{u}(\tau) = r \begin{pmatrix} \sin\theta(\tau)\cos\phi(\tau) \\ \sin\theta(\tau)\sin\phi(\tau) \\ \cos\theta(\tau) \end{pmatrix}
\end{equation}
where the elevation $\theta(\tau)$ and azimuth $\phi(\tau)$ are parametrized by the manifold coordinate $\tau \in [0, 2\pi]$:
\begin{align}
\theta(\tau) &= \theta_0 + \Delta\theta \sin(2\tau) \\
\phi(\tau) &= \Delta\phi \sin(\tau)
\end{align}
Here, $\theta_0$ is the mean elevation angle, and $\Delta\theta$, $\Delta\phi$ are the path amplitudes.

\subsubsection{Resolution of the Radial Distance}
Since the physical tether between $A$ and $K$ is assumed rigid with a constant length $\ell$, the coordinates must satisfy the algebraic distance constraint:
\begin{equation}
\|\vec{x}_K - \vec{x}_A\|^2 = \ell^2 \iff \|r\vec{u}(\tau) - \vec{x}_A(\alpha)\|^2 = \ell^2
\end{equation}
Expanding this relation yields a quadratic equation in $r$:
\begin{equation}
r^2 - 2 D(\alpha, \tau) r + a^2 - \ell^2 = 0
\label{eq:quadratic_r}
\end{equation}
where $D(\alpha, \tau)$ is the projection of the arm's tip onto the direction of the kite:
\begin{equation}
D(\alpha, \tau) = \vec{x}_A(\alpha) \cdot \vec{u}(\tau) = a \sin\theta(\tau) \cos(\alpha - \phi(\tau))
\end{equation}
Solving Eq.~\eqref{eq:quadratic_r} for the outer positive physical root yields:
\begin{equation}
r(\alpha, \tau) = D(\alpha, \tau) + \sqrt{D(\alpha, \tau)^2 + \ell^2 - a^2}
\end{equation}

\subsubsection{Generalized Coordinates}
Depending on the modeling abstraction level, two configurations are studied:
\begin{itemize}
    \item \textbf{Simplified 2D formulation (for standard paper model)}: The kite is strictly constrained to the lemniscate path and the tether is inextensible. The configuration is completely described by 2 degrees of freedom:
    \begin{equation}
    q = (\alpha, \tau)^T
    \end{equation}
    \item \textbf{Unconstrained 5D formulation (for high-fidelity models)}: We relax the algebraic constraints of rigid path tracking and rigid tether. The configuration is described by 5 generalized coordinates:
    \begin{equation}
    q_5 = (\alpha, \tau, R, \theta_\ell, \phi_\ell)^T
    \end{equation}
    where $R$ is the dynamic tether length, and $\theta_\ell$, $\phi_\ell$ represent the polar and azimuth angles of the tether vector $\vec{\ell} = \vec{x}_K - \vec{x}_A$ respectively.
\end{itemize}

\subsection{Kinematics}

The physical velocities of points $A$ and $K$ are mapped from the generalized coordinate velocities $\dot{q} = (\dot{\alpha}, \dot{\tau})^T$ via the system Jacobians:
\begin{equation}
\vec{v}_X = \dot{\vec{x}}_X = \frac{\partial \vec{x}_X}{\partial q} \dot{q}
\end{equation}
For the arm tip $A$, the velocity is:
\begin{equation}
\vec{v}_A = \frac{\partial \vec{x}_A}{\partial \alpha} \dot{\alpha} = \begin{pmatrix} -a\sin\alpha \\ a\cos\alpha \\ 0 \end{pmatrix} \dot{\alpha}
\end{equation}
For the kite $K$, the velocity is computed as:
\begin{equation}
\vec{v}_K = \frac{\partial \vec{x}_K}{\partial \alpha} \dot{\alpha} + \frac{\partial \vec{x}_K}{\partial \tau} \dot{\tau}
\end{equation}
where the partial derivatives are explicitly given by:
\begin{align}
\frac{\partial \vec{x}_K}{\partial \alpha} &= \frac{\partial r}{\partial \alpha} \vec{u}(\tau) \\
\frac{\partial \vec{x}_K}{\partial \tau} &= \frac{\partial r}{\partial \tau} \vec{u}(\tau) + r \frac{\partial \vec{u}}{\partial \tau}
\end{align}
Applying the chain rule to the quadratic solution for $r$ yields:
\begin{equation}
\frac{\partial r}{\partial \alpha} = \frac{r}{\sqrt{D^2 + \ell^2 - a^2}} \frac{\partial D}{\partial \alpha}, \quad \frac{\partial r}{\partial \tau} = \frac{r}{\sqrt{D^2 + \ell^2 - a^2}} \frac{\partial D}{\partial \tau}
\end{equation}
where:
\begin{align}
\frac{\partial D}{\partial \alpha} &= -a \sin\theta(\tau) \sin(\alpha - \phi(\tau)) \\
\frac{\partial D}{\partial \tau} &= a \cos\theta(\tau) \theta'(\tau) \cos(\alpha - \phi(\tau)) + a \sin\theta(\tau) \phi'(\tau) \sin(\alpha - \phi(\tau))
\end{align}
and the temporal derivatives of the path angles are:
\begin{equation}
\theta'(\tau) = 2 \Delta\theta \cos(2\tau), \quad \phi'(\tau) = \Delta\phi \cos(\tau)
\end{equation}

\subsection{Dynamics}

The system experiences gravitational, generator-resisting, and aerodynamic forces. Figure \ref{fig:dynamics_config} outlines the force distribution on the coupled arm-kite system.

\begin{figure}[htbp]
\centering
\begin{tikzpicture}[tdplot_main_coords, scale=2.2]

% ==== ORIGIN AND AXES ====
\coordinate (O) at (0,0,0);
\node[circle,fill,inner sep=1pt] at (O) {};
\node[below left] at (O) {O};

\draw[->] (0,0,0) -- (3,0,0) node[anchor=west]{$x$};
\draw[->] (0,0,0) -- (0,1.5,0) node[anchor=south]{$y$};
\draw[->] (0,0,0) -- (0,0,2) node[anchor=west]{$z$};

% Setup coordinates exactly as in Fig 1
\def\a{2}
\def\alphaA{20}
\pgfmathsetmacro{\xA}{\a*cos(\alphaA)}
\pgfmathsetmacro{\yA}{\a*sin(\alphaA)}
\coordinate (A) at (\xA, \yA, 0);
\node[circle,fill,inner sep=1pt] at (A) {};
\node[below] at (A) {A};

\def\l{2}
\def\thetaZero{70}
\def\DeltaTheta{5}
\def\DeltaPhi{30}
\def\tauK{90}

\pgfmathdeclarefunction{thetaFunc}{1}{\pgfmathparse{\thetaZero + \DeltaTheta * sin(2*#1)}}
\pgfmathdeclarefunction{phiFunc}{1}  {\pgfmathparse{\DeltaPhi * sin(#1)}}
\pgfmathdeclarefunction{xTauFunc}{1} {\pgfmathparse{sin(thetaFunc(#1)) * cos(phiFunc(#1))}}
\pgfmathdeclarefunction{yTauFunc}{1} {\pgfmathparse{sin(thetaFunc(#1)) * sin(phiFunc(#1))}}
\pgfmathdeclarefunction{zTauFunc}{1} {\pgfmathparse{cos(thetaFunc(#1))}}

\pgfmathsetmacro{\dotprod}{\xA * xTauFunc(\tauK) + \yA * yTauFunc(\tauK)}
\pgfmathsetmacro{\rK}{\dotprod + sqrt(\dotprod * \dotprod + \l^2 - \a^2)}

\pgfmathsetmacro{\xK}{\rK * xTauFunc(\tauK)}
\pgfmathsetmacro{\yK}{\rK * yTauFunc(\tauK)}
\pgfmathsetmacro{\zK}{\rK * zTauFunc(\tauK)}

\coordinate (K) at (\xK, \yK, \zK);

% Draw arm and tether
\draw[thick] (O) -- (A);
\draw (A) -- (K);
\node[circle,fill,inner sep=1pt] at (K) {};
\node[above] at (K) {K};

% ==== WIND SHEAR BACKGROUND ====
\foreach \z in {0.3, 0.7, 1.1, 1.5, 1.9} {
    \pgfmathsetmacro{\windlen}{0.45 * (\z)^0.16}
    \draw[->, cyan!80!black, thin] (-0.5, -0.4, \z) -- ++(\windlen, 0, 0);
}
\node[cyan!80!black, right] at (-0.5, -0.4, 1.9) {$\vec{v}_w(z)$};

% ==== FORCES AT THE ORIGIN (Resisting Torque) ====
\draw[->, purple, thick] (0.4, 0, 0) arc (0:-100:0.4) node[midway, below left] {$\Gamma_r$};

% ==== FORCES AT THE KITE ====
% Gravity
\draw[->, red!80!black, thick] (K) -- ++(0, 0, -0.55) node[below] {$\vec{F}_{g}$};

% Apparent Wind vector
\draw[->, cyan, thick] (K) -- ++(0.5, -0.1, -0.1) node[right] {$\vec{v}_a$};

% Aerodynamic Lift & Drag
\draw[->, orange, thick] (K) -- ++(0.35, -0.07, -0.07) node[below right] {$\vec{F}_D$};
\draw[->, green!60!black, thick] (K) -- ++(-0.1, -0.15, 0.45) node[above] {$\vec{F}_L$};

% Tether Tension
\draw[->, blue, thick] (K) -- ($(K)!0.35!(A)$) node[midway, left] {$\vec{T}$};

\end{tikzpicture}
\caption{Dynamic force environment acting on the oscillating arm and the kite, displaying the wind profile $\vec{v}_w(z)$, resisting generator torque $\Gamma_r$, gravity $\vec{F}_g$, tether tension $\vec{T}$, aerodynamic drag $\vec{F}_D$, and lift $\vec{F}_L$.}
\label{fig:dynamics_config}
\end{figure}

\subsubsection{Wind Profile}
To capture realistic wind shear variations with altitude $z$, we assume a standard power-law wind profile:
\begin{equation}
\vec{v}_w(z) = v_{w, \text{ref}} \left(\frac{z}{z_{\text{ref}}}\right)^\gamma \vec{e}_x
\end{equation}
where $v_{w, \text{ref}}$ is the wind velocity at a chosen reference altitude $z_{\text{ref}}$, and $\gamma$ is the atmospheric shear exponent.

\subsubsection{Aerodynamic Forces}
The relative velocity of the air with respect to the kite, known as the apparent wind velocity, is:
\begin{equation}
\vec{v}_a = \vec{v}_w(z_K) - \vec{v}_K
\end{equation}
The aerodynamic forces consist of lift $\vec{F}_L$ and drag $\vec{F}_D$:
\begin{align}
\vec{F}_D &= \frac{1}{2} \rho_a A_K C_D \|\vec{v}_a\|^2 \vec{e}_D \\
\vec{F}_L &= \frac{1}{2} \rho_a A_K C_L \|\vec{v}_a\|^2 \vec{e}_L
\end{align}
where:
\begin{itemize}
    \item $\rho_a$ is the air density, and $A_K$ is the reference wing area of the kite.
    \item $C_L$, $C_D$ are the lift and drag coefficients.
    \item $\vec{e}_D = \frac{\vec{v}_a}{\|\vec{v}_a\|}$ is the drag unit vector, pointing in the direction of the apparent wind flow.
    \item $\vec{e}_L$ is the lift unit vector, which is orthogonal to $\vec{e}_D$ and lies in the kite's symmetry plane.
\end{itemize}
The total aerodynamic force acting on the kite is the vector sum:
\begin{equation}
\vec{F}_{\text{aer}} = \vec{F}_L + \vec{F}_D
\end{equation}
Tether drag can also be lumped directly into the kite's drag coefficient $C_D$ by adding an equivalent drag contribution proportional to the tether dimensions.

\subsubsection{Gravity}
Gravity acts vertically downwards on both the kite (mass $m_K$) and the arm (distributed mass or concentrated mass):
\begin{equation}
\vec{F}_g = -m_K g \vec{e}_z
\end{equation}

\subsubsection{Resisting Torque}
An electric generator connected to the arm pivot at $O$ extracts electrical energy by exerting a damping/control resisting torque $\Gamma_r$ that opposes the rotation of the arm:
\begin{equation}
\Gamma_r = -c_{\text{gen}} \dot{\alpha}
\end{equation}
where $c_{\text{gen}} > 0$ is the generator's damping/load coefficient.

\subsubsection{Lagrangian Mechanics Formulation}
Applying the Euler-Lagrange equations to the generalized coordinates $q = (\alpha, \tau)^T$, the equations of motion are expressed as:
\begin{equation}
M(q)\ddot{q} + C(q,\dot{q})\dot{q} + G(q) = Q(q,\dot{q})
\end{equation}
where $M(q)$ is the system mass matrix, $C(q,\dot{q})$ captures the Coriolis and centrifugal terms, $G(q)$ represents the generalized gravity forces, and $Q(q,\dot{q})$ contains the projections of the non-conservative forces (aerodynamic forces $\vec{F}_{\text{aer}}$ and resisting generator torque $\Gamma_r$) onto the system tangent space.

\section*{References}
\noindent [1] A. Bavoil, J.-B. Caillau, L. Dell'Elce, and A. N\^eme. "Airborne wind energy: the KEEP project." Accepted at the European Control Conference (ECC), 2026.

\end{document}